%=====================================================================
% FRONT MATTER
%=====================================================================
\pagenumbering{roman}
\structuralfigures{F.}
\thispagestyle{empty}

\begin{center}
\vspace*{1.0cm}

{\color{secondary}\rule{0.85\textwidth}{2pt}}\\[6pt]
{\large\textsc{\color{secondary}Bachelor of Science in Information Technology}}\\[18pt]

{\Huge\bfseries\color{primary}E-COMMERCE APPLICATION\\[4pt]DEVELOPMENT}\\[10pt]
{\LARGE\bfseries\color{primary}Building, Running and Securing an Online Shop}\\[16pt]

{\color{secondary}\rule{0.85\textwidth}{2pt}}\\[16pt]

{\Large\textsc{Unified Lecture Handout}}\\[6pt]
{\large\color{gray!70!black}Complete Course in One Volume $\bullet$ Fall 2026}\\[24pt]

% --- the claim the whole book rests on ---
\begin{tikzpicture}
  \node[rounded corners=4pt, draw=primary, line width=1.1pt, fill=primary!6,
        text=primary, align=center, font=\small\bfseries,
        minimum width=11.6cm, minimum height=1.5cm]
       {\faIcon{store}\\[2pt]Anyone can install a storefront in an afternoon.
        This course is everything after that.};
\end{tikzpicture}\\[10pt]
{\small\itshape\color{gray!70!black}A shop is not a website. It takes money,
holds other people's data, and is answerable in law --- and each of those
changes what ``finished'' means.}

\vspace{1.0cm}

\begin{tcolorbox}[enhanced, width=0.86\textwidth, colback=lightbg,
  colframe=primary, boxrule=0pt, leftrule=4pt, arc=2pt, top=6pt, bottom=6pt]
\small
\textbf{\color{primary}What this volume covers.} Every distinct clause of the
HEC outline for \emph{E-Commerce}, in the order the outline gives them, traced
chapter by chapter in the Course Specification overleaf.

\smallskip
That outline is published once only, in the \textbf{2017} curriculum (p.~94),
where the course is an \emph{elective in BS Software Engineering} rather than a
BSIT course at all. The 2023 curriculum drops it completely; the 2025
curriculum renames it \emph{Digital Marketing \& E-Commerce} and publishes no
contents. The department teaches it as \textbf{INFT-4041}, BSIT, 5th semester,
3 credit hours. All of that is recorded in the Course Specification rather
than smoothed over.

\smallskip
The course is taught on a \textbf{running shop}. One storefront is installed in
Chapter~5 and every chapter afterwards adds a layer to it; by Chapter~23 it is
the semester project. Appendix~A sets up the stack in one sitting.
\end{tcolorbox}

\vspace{0.9cm}
{\small\color{gray!60!black}Six Parts $\bullet$ Twenty-Three Chapters $\bullet$ Seven Appendices}
\vspace{1cm}
\end{center}

\newpage

%=====================================================================
\chapter*{How to Use This Handout}
\addcontentsline{toc}{chapter}{How to Use This Handout}
\markboth{How to Use This Handout}{}

There are two ways to teach this subject badly.

The first is to teach it as business studies: models, channels, trends, the
history of the dot-com boom, and a graduate who can describe an online shop but
has never configured a tax rule. The second is to teach it as programming
exercises: write a cart class, write a checkout function, and a graduate who has
built a toy that no one would take money through.

This handout takes the third path, which is also the one the published outline
describes if you read it closely. You install a real, open-source storefront in
Chapter~5, and you spend the rest of the course making it into a shop that could
actually trade: a catalogue that models real products, a checkout that handles
shipping and tax, a payment path that survives a failed transaction, an admin
who can process a refund, a server that does not leak its database, and a listing
a search engine can find. Where the platform cannot do what you need --- and in
Chapter~19 it cannot --- you write the code that extends it.

\section*{The six parts}

\textbf{Part~I --- E-Commerce and Its Models} (Chapters 1--4) is the business
grounding, and it is short on purpose. What electronic commerce actually changed;
the models and where the revenue comes from; auctions, portals and communities;
and the build-buy-or-extend decision that justifies everything the rest of the
book does.

\textbf{Part~II --- Standing the Shop Up} (Chapters 5--8) gets a storefront
running and fills it. The stack and what each layer does; products and categories,
with the tables underneath them; variants, options and uploads; then the theme.

\textbf{Part~III --- From Basket to Order} (Chapters 9--12) is the pipeline a
customer walks through. Where basket state lives; the checkout and the order
state machine; shipping, tax and zones; discounts, vouchers and referrals.

\textbf{Part~IV --- Money and Identity} (Chapters 13--16) is the part where
mistakes cost money. Payment systems and who carries the risk; integrating a
gateway and testing it in a sandbox; accounts, sessions and password storage;
then securing the whole thing.

\textbf{Part~V --- Running and Extending the Shop} (Chapters 17--20) is the half
of the job nobody demonstrates: the dashboard and what to watch, the back office,
writing your first module, and deployment with everything that follows it.

\textbf{Part~VI --- Reach, Law and Practice} (Chapters 21--23) closes with being
found, being lawful, and your project.

\begin{conceptbox}
Why Chapter~19 is the hinge of the course. Everything before it is
\emph{configuration}: you are choosing among the options a platform already
offers, and a careful non-programmer could do most of it. Chapter~19 is the
first time the answer to ``can the shop do this?'' is no, and you have to write
the code that makes it yes. That is the difference between the outline's bare
title, \emph{E-Commerce}, and the one this department teaches,
\emph{E-Commerce Application Development}.
\end{conceptbox}

\section*{The furniture of every chapter}

\rowcolors{2}{white}{lightbg}
\begin{longtable}{@{}L{4.6cm}L{10.6cm}@{}}
\toprule
\rowcolor{primary!15}
\textbf{Box} & \textbf{What it is for} \\
\midrule
\endhead
\faIcon{bullseye}~\textbf{Outcomes} & What you will be able to \emph{do}. Bloom
verbs, not topics. \\
\faIcon{link}~\textbf{Before you start} & Which earlier chapters this one stands
on, and what state your shop should be in. \\
\faIcon{tags}~\textbf{Key terms} & The vocabulary introduced; all of it is in the
glossary of Appendix~D. \\
\faIcon{book}~\textbf{Definition} & The canonical statement. In this subject
several ordinary English words --- order, basket, customer, product --- are
technical terms with narrower meanings, and confusing the two causes real bugs. \\
\faIcon{lightbulb}~\textbf{Concept} & The judgement call: what an experienced
practitioner would decide, and why. \\
\faIcon{exclamation-triangle}~\textbf{Warning} & Something that will bite you,
usually in production rather than in the examination. \\
\faIcon{calculator}~\textbf{Worked example} & Real arithmetic --- a tax
calculation, a gateway fee, a discount applied in the wrong order --- with every
step shown. \\
\faIcon{tachometer-alt}~\textbf{Measured} & A number produced by running
something, not by asserting it. Twelve chapters carry one; see below. \\
\faIcon{exclamation-circle}~\textbf{Common pitfalls} & The mistakes that are
actually made. \\
\faIcon{flask}~\textbf{Try it yourself} & The laboratory. Every one of them
changes the same shop, and they are cumulative: skip Chapter~9's and
Chapter~10's will have nothing to check out. \\
\faIcon{question-circle}~\textbf{Checkpoint} & Three questions. Attempt them
before reading Appendix~C. \\
\faIcon{clipboard-check}~\textbf{Chapter summary} & What to retain. \\
\faIcon{pen-fancy}~\textbf{Review questions} & Examination practice: multiple
choice, short answer, applied. \\
\bottomrule
\end{longtable}

\newpage

Two typographic conventions carry a lot of weight in a course taught through a
graphical interface, and they are worth learning on this page:

\begin{itemize}[leftmargin=1.6em, itemsep=3pt, topsep=4pt]
  \item \adminpath{Catalog, Products, Add New} is a \textbf{click path} through
        the administration panel. Read the chevrons as ``then''.
  \item \uiel{Save} on its own is a single control --- a button, a field, a tab
        or a menu item. Anything in \texttt{typewriter} is text you type or a
        file you edit, which is a different thing entirely.
\end{itemize}

\begin{alertbox}[The measured strand, and why it is only in twelve chapters]
Twelve of the twenty-three chapters end with a number that was produced rather
than asserted: a page time, a query count, a request rate, a hash cost, a list
of security headers present and absent.

\smallskip
They are the twelve where measuring genuinely teaches something. ``Add an index
and the category page gets faster'' is a claim a student can be told and will
forget; watching the same page go from a table scan at fifty thousand products
to an index seek, and reading both numbers off your own machine, is not. The
other eleven chapters argue in prose, because dressing up a legal obligation or
a business model with a stopwatch would be theatre.

\smallskip
Every figure in a \faIcon{tachometer-alt}~\textbf{Measured} box was produced by
the file named beside it, in \texttt{code/}, against the stack in Appendix~A.
Rerun them. Your numbers will differ from the book's --- different machine,
different disk, different day --- and the \emph{shape} should not. Where a
result depends on the particular machine, the box says so.
\end{alertbox}

\newpage

%=====================================================================
\chapter*{Course Specification}
\addcontentsline{toc}{chapter}{Course Specification}
\markboth{Course Specification}{}

This page is the course file. It states what HEC requires of a course called
\emph{E-Commerce} and shows where each requirement is met.

That is harder here than in most courses, because the requirement has been
published exactly once and then withdrawn. The three editions are traced below
so that an accreditation review against any of them has something to check.

\section*{Course particulars}

\rowcolors{2}{white}{lightbg}
\begin{longtable}{@{}L{4.2cm}L{11.0cm}@{}}
\toprule
\rowcolor{primary!15}
\textbf{Item} & \textbf{Value} \\
\midrule
\endhead
Course title & \emph{E-Commerce Application Development}, as taught. HEC 2017
calls it simply \emph{E-Commerce}; HEC 2025 calls it \emph{Digital Marketing \&
E-Commerce}. The department's title is the accurate one, and the Course
Specification justifies the two added words \\
Course code & INFT-4041 \\
Programme & BS Information Technology, 5th semester. Under HEC 2017 the course
is an \textbf{elective in BS Software Engineering} (2017, p.~38, ``select any
FIVE'') and does not appear in the BSIT scheme at all \\
Credit hours & 3: three credit hours of lectures, no separate laboratory credit.
The laboratory work in this handout is therefore done in lecture time and as
homework, which is why every lab is small and cumulative \\
Prerequisite & \emph{Web Technologies} (CS-207 / INFT-4005, 4th semester). HEC
2017 names the prerequisite \emph{Web Engineering}; the department's course of
that content is called Web Technologies. The name differs, the dependency is
satisfied --- see Appendix~F \\
Duration & Sixteen teaching weeks, midterm in Week~8, project due Week~16
(Appendix~G) \\
Teaching methodology & As published (HEC 2017): lecturing, written assignments,
\textbf{project}, report writing \\
Assessment & As published: sessional examination, home assignments, quizzes,
project, presentations, final examination \\
Set texts & As published: Laudon and Traver, \emph{E-Commerce}; Peacock,
\emph{PHP 5 E-commerce Development}; Rayport, \emph{Introduction to E-Commerce};
Schneider, \emph{Electronic Commerce}. See defect~5 below, and Appendix~E \\
\bottomrule
\end{longtable}

\section*{The three editions}

\rowcolors{2}{white}{lightbg}
\begin{longtable}{@{}L{1.5cm}L{3.6cm}L{6.4cm}L{3.3cm}@{}}
\toprule
\rowcolor{primary!15}
\textbf{Ed.} & \textbf{Title} & \textbf{Status} & \textbf{Contents} \\
\midrule
\endhead
2017 & E-Commerce & Elective in BS Software Engineering (p.~38); absent from
BSIT. 3(3,0); prerequisite \emph{Web Engineering} & \textbf{Yes, p.~94 --- the
only published outline in any edition} \\
2023 & --- & \textbf{Absent entirely.} The course appears nowhere in the
computing-disciplines curriculum except as an incidental application domain
(p.~23) & No \\
2025 & Digital Marketing \& E-Commerce & Interdisciplinary / Allied course,
item~11 (p.~16). Not a computing core; universities may add others & No --- title
and credit hours only \\
\bottomrule
\end{longtable}

\begin{alertbox}[Five defects in the published document]
Recorded here so that a reader comparing this handout with the HEC booklet is
not confused, and so that the department's own course file does not inherit
them.

\smallskip
\textbf{1. Two learning outcomes, both too broad to assess.} HEC publishes
exactly two: ``Understand the concepts and standards related to the discipline
of E-Commerce'' and ``Analyze complex real world problems found in E-Commerce''.
Sibling courses publish eight. Two cannot carry sixteen weeks or a question
paper, so this department adds six, each traced back to one of HEC's two and
each marked as departmental in the table below.

\smallskip
\textbf{2. The Bloom's Taxonomy column is blank} for both outcomes --- only the
Domain column (``C'') is filled. The levels below are therefore
\emph{assigned by this department}, and are marked as such.

\smallskip
\textbf{3. The ``Course Content'' paragraph is the table of contents of its own
reference book.} Clause for clause, in order, it is the chapter list of
Michael Peacock's \emph{PHP 5 E-commerce Development} (Packt, 2010), which is
reference~2 on the same page. HEC did not write an outline for this course; it
pasted a contents page.

\smallskip
\textbf{4. Two clauses are duplicated}, which is what pasting a contents page
does. ``Checkout'' appears twice --- once as ``The Checkout and Order Process''
--- and ``Managing Products and Categories'' appears twice, once on its own and
once under Administration. Twenty-four clauses; twenty-two distinct.

\smallskip
\textbf{5. The prescribed technology is end-of-life and the editions are stale.}
PHP~5 lost security support on 31~December 2018: the prescribed text targets a
runtime that has been unsafe to deploy for longer than most of this class has
been studying. Laudon and Traver is cited at the 13th edition (2017) and Rayport
at 2007. Appendix~E cites current editions and explains each gap.
\end{alertbox}

\section*{Coverage of the HEC 2017 course outline}

The published outline is a single run-on paragraph. Each of its clauses is
traced below, in the order HEC writes them. Nothing is omitted, and the two
duplicates are shown where they fall rather than quietly merged.

\rowcolors{2}{white}{lightbg}
\begin{longtable}{@{}L{6.4cm}L{2.0cm}L{6.6cm}@{}}
\toprule
\rowcolor{primary!15}
\textbf{HEC 2017 clause} & \textbf{Chapter} & \textbf{Where exactly} \\
\midrule
\endhead
``An overview of E-Commerce \& its business models and concepts'' & 1, 2 & What
changed and why; then the models, with the revenue model treated as the
question that decides the rest \\
``Planning an E-Commerce Framework'' & 4 & Requirements, and the
build-buy-or-extend decision this course's answer rests on \\
``Managing Products and Categories'' & 6 & The admin screens, and the tables
beneath them \\
``Product Variations and User Uploads'' & 7 & Options and variants, the
combinatorics they generate, and why an upload field is a security question \\
``Enhancing the User Experience'' & 8 & Themes and template overrides, measured
by page weight rather than taste \\
``The Shopping Basket'' & 9 & Session, cookie or database, and what survives a
restart \\
``The Checkout and Order Process'' & 10 & The order state machine, step by
step \\
``Shipping and Tax'' & 11 & Geo zones, tax classes, and sales tax on a
Pakistani online sale \\
``Discounts, Vouchers, and Referrals'' & 12 & All three, and the order in which
they must be applied \\
``Checkout'' \emph{(duplicate of clause~7)} & 10 & Same chapter; see defect~4 \\
``Taking Payment for Orders'' & 14 & The gateway handshake, the sandbox, and
what happens when it fails halfway \\
``User Account Management'' & 15 & Registration, sessions, password storage,
reset flows \\
``Administration: Dashboard'' & 17 & The dashboard, and which of its numbers
are worth acting on \\
``Managing Products and Categories'' \emph{(duplicate of clause~3)} & 6, 17 &
Same material from the back office; see defect~4 \\
``Managing Orders'' & 18 & Order states, fulfilment, partial shipment \\
``Customers'' & 18 & Customer records, groups and approval \\
``Refunds'' & 18 & Returns, credit and what the order state machine does next \\
``Voucher Codes'' & 12, 18 & Created in 12, administered in 18 \\
``Shipping'' \emph{(administration)} & 11, 18 & Configured in 11, operated
in 18 \\
``Deploying, Security, and Maintenance'' & 16, 20 & Security is given its own
chapter --- see the note below --- and deployment, backup and maintenance
share 20 \\
``Web Payment Systems'' & 13 & Cards, wallets, bank transfer and cash on
delivery; cost and risk for each \\
``Social, Legal, and Ethical Issues of E-Commerce'' & 22 & Consumer rights,
privacy, dark patterns, accessibility, and the Pakistani statute \\
``Auctions, Portals, and Communities'' & 3 & Moved forward to Part~I, where it
belongs: these are business models \\
``SEO'' & 21 & Crawlability, URL structure, metadata, sitemaps \\
\bottomrule
\end{longtable}

\section*{Course learning outcomes}

The two outcomes HEC publishes, verbatim, followed by the six this department
adds. Bloom levels are assigned by this department throughout, because the
published table leaves the column blank.

\rowcolors{2}{white}{lightbg}
\begin{longtable}{@{}L{1.3cm}L{7.6cm}L{2.1cm}L{3.8cm}@{}}
\toprule
\rowcolor{primary!15}
\textbf{CLO} & \textbf{By the end of the course the student will be able to} &
\textbf{Bloom} & \textbf{Chapters} \\
\midrule
\endhead
CLO-1 & \emph{(HEC, verbatim)} Understand the concepts and standards related to
the discipline of E-Commerce & C2 Understand & 1--4, 13, 22 \\
CLO-2 & \emph{(HEC, verbatim)} Analyze complex real world problems found in
E-Commerce & C4 Analyse & 6, 10, 16, 20, 23 \\
CLO-3 & \emph{(dept.)} Explain e-commerce business and revenue models, and
classify a given venture against them & C2 Understand & 2, 3 \\
CLO-4 & \emph{(dept.)} Install, configure and secure an open-source storefront
on a reproducible stack & C3 Apply & 5, 16, 20 \\
CLO-5 & \emph{(dept.)} Model a catalogue --- products, categories, options and
variants --- and justify the schema that results & C3 Apply & 6, 7 \\
CLO-6 & \emph{(dept.)} Configure the order pipeline end to end: basket,
checkout, shipping, tax, discounts and payment & C3 Apply & 9--12, 14 \\
CLO-7 & \emph{(dept.)} Analyse a storefront's performance and security posture
from measurement rather than assertion & C4 Analyse & 6, 8, 15, 16, 20, 21 \\
CLO-8 & \emph{(dept.)} Extend the platform with a module meeting a stated
requirement, then deploy, present and defend the result & C5 Evaluate &
19, 20, 23 \\
\bottomrule
\end{longtable}

\begin{conceptbox}[Three topics taught that the outline does not name]
Each is justified here rather than smuggled in, because a reviewer will ask.

\smallskip
\textbf{Chapter~6's data model.} The outline names the screens --- ``Managing
Products and Categories'' --- but never the tables. A graduate who can fill in
an admin form but cannot say why a product's name lives in a different table
from its price has not understood the catalogue, and \emph{Database Systems} is
a prerequisite that ought to be paid off rather than left on the shelf.

\smallskip
\textbf{Chapter~16, Security, as a chapter of its own.} HEC buries it inside
``Deploying, Security, and Maintenance'', one clause among three. A shop holds
addresses, order histories and password hashes, and CLO-7 asks for a security
posture to be \emph{analysed}. One third of one clause will not do it.

\smallskip
\textbf{Chapter~19, module development.} The outline is configuration from end
to end; nothing in it requires a student to write a line of code. The
department's title says \emph{Application Development}, the HEC teaching
methodology requires a project, and a project that is entirely configuration is
not a computing project.
\end{conceptbox}

\begin{alertbox}[Where this course sits, and the one awkwardness]
The outline this handout traces was written for \emph{software engineering}
students as an elective, and it shows: it assumes the reader will build a cart
from scratch in PHP. This department teaches it to \emph{information technology}
students as a 5th-semester course, where the useful skill is standing up,
configuring, securing and extending a real platform rather than reimplementing
one badly.

\smallskip
This handout therefore follows the outline's \emph{clauses} faithfully and its
\emph{method} not at all. Every topic HEC lists is taught; none of it is taught
by writing a shopping cart in PHP~5. Chapter~4 makes that argument explicitly,
so that the student can defend the choice rather than merely inherit it.
\end{alertbox}

\newpage

%=====================================================================
\chapter*{The Course at a Glance}
\addcontentsline{toc}{chapter}{The Course at a Glance}
\markboth{The Course at a Glance}{}

\dsfigh{%
\begin{tikzpicture}[
  part/.style={rounded corners=3pt, fill=#1, text=white, font=\bfseries\footnotesize,
               minimum width=2.5cm, text width=2.2cm, minimum height=0.95cm,
               align=center},
  ch/.style={rounded corners=2pt, draw=#1, line width=0.7pt, fill=#1!7, text=#1!60!black,
             font=\scriptsize, minimum width=2.5cm, minimum height=0.52cm, align=center},
  arr/.style={-{Stealth[length=2mm]}, primary!45, line width=0.8pt}
]
\node[part=primary]   (p1) at (0,0)    {Part I\\Models};
\node[part=secondary] (p2) at (3.0,0)  {Part II\\Standing It Up};
\node[part=greenhl]   (p3) at (6.0,0)  {Part III\\Basket to Order};
\node[part=purplehl]  (p4) at (9.0,0)  {Part IV\\Money and Identity};
\node[part=tealhl]    (p5) at (12.0,0) {Part V\\Running It};
\node[part=goldhl]    (p6) at (15.0,0) {Part VI\\Reach and Law};

\foreach \a/\b in {p1/p2,p2/p3,p3/p4,p4/p5,p5/p6}{\draw[arr] (\a) -- (\b);}

\foreach \i/\t in {1/{1. What changed}, 2/{2. Business models},
                   3/{3. Auctions, portals}, 4/{4. Build or extend}}{
  \pgfmathsetmacro{\yy}{-0.90-(\i-1)*0.62}
  \node[ch=primary] at (0,\yy) {\t};}

\foreach \i/\t in {1/{5. The stack}, 2/{6. Catalogue},
                   3/{7. Variations}, 4/{8. Theme and UX}}{
  \pgfmathsetmacro{\yy}{-0.90-(\i-1)*0.62}
  \node[ch=secondary] at (3.0,\yy) {\t};}

\foreach \i/\t in {1/{9. The basket}, 2/{10. Checkout},
                   3/{11. Shipping, tax}, 4/{12. Discounts}}{
  \pgfmathsetmacro{\yy}{-0.90-(\i-1)*0.62}
  \node[ch=greenhl] at (6.0,\yy) {\t};}

\foreach \i/\t in {1/{13. Payment systems}, 2/{14. The gateway},
                   3/{15. Accounts}, 4/{16. Security}}{
  \pgfmathsetmacro{\yy}{-0.90-(\i-1)*0.62}
  \node[ch=purplehl] at (9.0,\yy) {\t};}

\foreach \i/\t in {1/{17. Dashboard}, 2/{18. Orders, refunds},
                   3/{19. Your first module}, 4/{20. Deployment}}{
  \pgfmathsetmacro{\yy}{-0.90-(\i-1)*0.62}
  \node[ch=tealhl] at (12.0,\yy) {\t};}

\foreach \i/\t in {1/{21. SEO}, 2/{22. Law and ethics},
                   3/{23. Project}}{
  \pgfmathsetmacro{\yy}{-0.90-(\i-1)*0.62}
  \node[ch=goldhl] at (15.0,\yy) {\t};}

\node[rounded corners=3pt, fill=primary!7, draw=primary!30,
      minimum width=17.3cm, minimum height=0.66cm,
      font=\scriptsize\bfseries, text=primary] at (7.5,-4.35)
      {\faIcon{store}~~One shop, installed in Chapter 5 and added to every week
       until it is the project in Chapter 23.};
\end{tikzpicture}}%
{The six parts and the chapters in each. Chapter 19 is the hinge: it is the
first chapter in which configuring the platform is not enough and you have to
write code.}{coursemap}

\vspace{4pt}
{\small\itshape\color{gray!70!black}
This course stands on \emph{Web Technologies} and \emph{Database Systems}, and
shares its deployment ground with \emph{Virtual Systems and Services} in the
same semester. Appendix~F states what belongs to each.}

\newpage

%=====================================================================
\chapter*{The Shop We Build}
\addcontentsline{toc}{chapter}{The Shop We Build}
\markboth{The Shop We Build}{}

Every lab in this book changes the same shop. Nothing is thrown away, and
nothing is built twice. This page is the map of what gets added when, so that
you can tell at any point in the semester whether you are behind.

\dsfigh{%
\begin{tikzpicture}[
  lay/.style={rounded corners=2pt, draw=#1, line width=0.8pt, fill=#1!8,
              text=#1!55!black, font=\scriptsize, minimum height=0.62cm,
              text width=4.6cm, align=center},
  sidenote/.style={font=\scriptsize\bfseries, text=primary, anchor=west}
]
\node[rounded corners=3pt, fill=primary, text=white, font=\bfseries\footnotesize,
      minimum width=5.0cm, minimum height=0.8cm, align=center] (base) at (0,0)
     {Ch 5\\The stack, running};

\foreach \i/\c/\t in {
   1/secondary/{Ch 6--7: a catalogue},
   2/secondary/{Ch 8: a theme},
   3/greenhl/{Ch 9--10: basket and checkout},
   4/greenhl/{Ch 11--12: shipping, tax, vouchers},
   5/purplehl/{Ch 14: a payment gateway},
   6/purplehl/{Ch 15: customer accounts},
   7/purplehl/{Ch 16: hardened},
   8/tealhl/{Ch 19: your own module},
   9/goldhl/{Ch 20--21: deployed, indexed}}{
  \pgfmathsetmacro{\yy}{0.95+(\i-1)*0.78}
  \node[lay=\c] at (0,\yy) {\t};}

\draw[-{Stealth[length=2mm]}, primary!40, line width=1pt]
      (-3.3,-0.45) -- (-3.3,8.15);
\node[rotate=90, font=\scriptsize\bfseries, text=primary!70]
      at (-3.65,3.85) {sixteen weeks};

\node[rounded corners=3pt, fill=goldhl, text=white, font=\bfseries\footnotesize,
      minimum width=5.0cm, minimum height=0.8cm, align=center] at (0,8.55)
     {Ch 23\\Defended as the project};

% Each note sits at the mid-height of the layers it describes, which is why
% these y-values are not evenly spaced: the layers are, the groups are not.
\node[sidenote, text width=6.0cm] at (2.9,0)    {Docker Compose brings up the
   storefront, the database and a database console. Nothing is installed on
   your own machine.};
\node[sidenote, text width=6.0cm] at (2.9,1.34) {Real products, with real option
   combinations, in a schema you can explain --- and a shop that looks like
   somebody's.};
\node[sidenote, text width=6.0cm] at (2.9,2.90) {A customer can now get from the
   front page to a confirmed order, at the right price.};
\node[sidenote, text width=6.0cm] at (2.9,4.85) {That order is now paid for, by
   an identified customer, over a shop that has been attacked and hardened.};
\node[sidenote, text width=6.0cm] at (2.9,6.41) {The platform does something it
   was not shipped able to do.};
\node[sidenote, text width=6.0cm] at (2.9,7.19) {Somebody who is not you can
   reach it, and find it.};
\end{tikzpicture}}%
{What each part adds to the one shop. A lab skipped is a lab the next chapter
cannot build on --- Chapter 10 has nothing to check out if Chapter 9's basket
was never configured.}{shoplayers}

\begin{alertbox}[Before Week 2]
The stack runs in containers, so nothing in this course asks you to install a
web server, PHP or a database on your own machine. It does ask you to install
\textbf{one} thing: a container runtime. Appendix~A does it in one sitting and
lists what to do when it will not start, which on a university laptop it
sometimes will not.

\smallskip
Do it before Week~2. From Chapter~5 onward every chapter assumes a shop you can
reach in a browser, and a student without one spends the semester reading about
a course the rest of the class is taking.
\end{alertbox}

\newpage

%=====================================================================
\chapter*{Prerequisite Self-Check}
\addcontentsline{toc}{chapter}{Prerequisite Self-Check}
\markboth{Prerequisite Self-Check}{}

Twelve questions. If you can answer ten, begin at Chapter~1. If you cannot,
Appendix~B covers the PHP, SQL and template syntax this book assumes, your
\emph{Web Technologies} and \emph{Database Systems} notes cover the rest --- and
the gap is worth closing in Week~1 rather than Week~9.

\begin{enumerate}[leftmargin=2.4em, itemsep=4pt, topsep=4pt]
  \item What is the difference between an HTTP \texttt{GET} and an HTTP
        \texttt{POST}, and which one must a form that charges a card use? Why?
        \hfill{\scriptsize\color{neutral}(Web Technologies)}
  \item What do the status codes 200, 302, 404 and 500 mean, and which one does
        a browser follow without telling the user?
        \hfill{\scriptsize\color{neutral}(Web Technologies)}
  \item HTTP is stateless. How, then, does a server know that two requests came
        from the same visitor?
        \hfill{\scriptsize\color{neutral}(Web Technologies, \chref{basket})}
  \item Write a \texttt{SELECT} that joins two tables and returns one row per
        match. \hfill{\scriptsize\color{neutral}(Database Systems, Appendix~B)}
  \item What is a database index for, and what does adding one cost?
        \hfill{\scriptsize\color{neutral}(Database Systems, \chref{catalogue})}
  \item Distinguish a primary key from a foreign key, and say what the database
        guarantees about each.
        \hfill{\scriptsize\color{neutral}(Database Systems)}
  \item One product has many photographs. Which table holds the photographs, and
        why not the product table?
        \hfill{\scriptsize\color{neutral}(Database Systems, \chref{catalogue})}
  \item In an HTML form, what does the \texttt{name} attribute do, and where
        does its value arrive on the server?
        \hfill{\scriptsize\color{neutral}(Web Technologies)}
  \item A theme sets a heading to 18\,px and you want 24\,px. Give two ways to
        override it in CSS, and say which one you would not use.
        \hfill{\scriptsize\color{neutral}(Web Technologies, \chref{ux})}
  \item In PHP, what does \texttt{\$\_POST['qty']} contain, and why may you
        never use it in a query without doing something to it first?
        \hfill{\scriptsize\color{neutral}(Appendix~B, \chref{security})}
  \item What does HTTPS protect, and name one thing it does \emph{not} protect.
        \hfill{\scriptsize\color{neutral}(Information Security, \chref{security})}
  \item Open a terminal in a named directory, run a command, and read its
        output. If that sentence is not routine, start with Appendix~A.
        \hfill{\scriptsize\color{neutral}(Appendix~A)}
\end{enumerate}

\begin{conceptbox}
Notice what is \emph{not} on that list. You do not need to be a strong PHP
programmer: until Chapter~19 you will read far more code than you write, and
Appendix~B gives you enough syntax to read it. You do not need any business
background. You do not need a powerful machine.

\smallskip
What you do need is the habit of \emph{checking} --- opening the database to
see what the form actually wrote, reading the error log rather than guessing,
rerunning a measurement after changing one thing. Most of what goes wrong in
this course is invisible from the front page of the shop, and every chapter
tries to teach you where to look.
\end{conceptbox}

\newpage
\tableofcontents
\newpage
\listoffigures
\newpage

\pagenumbering{arabic}
